\chapter{Esperimenti e risultati}
\label{ch:risultati}

% dataset (prendere 20 campioni, arrivare a 100 sarebbe l'ideale) https://github.com/huangd1999/UnLeakedTestBench/blob/main/datasets/ULT_Lite.jsonl

% modelli da https://build.nvidia.com/models?filters=usecase%3Ausecase_code_gen (prendere llama 1, 3, 8 billion)
% settare la temperatura a zero

% dichiarato il template per il prompt
% formattare il prompt
% [instruction] write the unit test for this function. Output only the code, formatted as ....
% [data] <sample del dataset>
%
% output viene salvato in file separati

% metriche da misurare
% - quanto sono in grado di coprire il codice effettivamente (coverage)
% - possibile misura di quanto sono rieseguite le linee (contare quante volte viene eseguita la stessa linea)
% - quante linee sono replicate (limitato ai file di test, per esempio con questo tool https://github.com/platisd/duplicate-code-detection-tool)
% codice non eseguibile -> errori di python o simili OPPURE test falliti (sarebbe l'ideale distinguerli)
% - eseguibile ma fallito
% - non eseguibile

% aggiunge ulteriore valore fare la valutazione "senza LLM", ovvero con AST + KLARA

% si può mettere anche un esempio di generazione di test end-to-end e la spiegazione delle metriche fatta per via visiva


% --- da qui: struttura da riempire con i risultati definitivi ---

\section{Setup sperimentale}
% Riepilogo compatto: 100 campioni di ULT_Lite \cite{huang2025ult}, tre modelli
% \cite{llama3herd} di taglia 1B, 3B e 8B, temperatura 0, limite di 1024 token,
% una funzione per richiesta, un file di output per coppia (modello, funzione).
% Dichiarare l'infrastruttura di esecuzione e la quantizzazione.


\section{Metriche misurate}
% Definizione operativa di ciascuna metrica, con la formula quando serve:
% - esito: non eseguibile / eseguibile ma fallito / passato;
% - copertura di istruzioni sulla funzione sotto test, con i due denominatori
%   (soli file eseguibili, oppure tutti i campioni);
% - proporzione di test corretti sul totale dei test generati;
% - riesecuzione delle righe: quante volte la stessa riga viene eseguita;
% - duplicazione fra i file di test;
% - aderenza al formato e tasso di troncamento.


\section{Risultati per modello}
% Tabella principale: una riga per modello, colonne = metriche.
% Grafici suggeriti:
% - barre impilate degli esiti per modello;
% - distribuzione della copertura (box plot) per modello;
% - dispersione copertura contro proporzione di test corretti.

\begin{table}[H]
\centering
\caption{Esiti dell'esecuzione dei test generati, per modello (100 campioni ciascuno).}
\label{tab:esiti}
\begin{tabular}{lcccc}
\hline
\textbf{Modello} & \textbf{Passati} & \textbf{Falliti} & \textbf{Non eseguibili} & \textbf{Copertura media} \\
\hline
Llama 3.2 1B & -- & -- & -- & -- \\
Llama 3.2 3B & -- & -- & -- & -- \\
Llama 3.1 8B & -- & -- & -- & -- \\
\hline
\end{tabular}
\end{table}


\section{Copertura e correttezza}
% Il risultato piu' rilevante emerso finora: la copertura resta alta mentre la
% proporzione di test corretti e' bassa. Il modello sceglie input plausibili e
% li fa arrivare al codice, ma sbaglia il valore atteso nell'assert.
% Osservazione importante: la copertura e' alta anche grazie ai test che
% falliscono, perche' un assert sbagliato esegue comunque la funzione prima di
% fallire. Le due metriche vanno lette insieme.
% Collegare a \cite{inozemtseva2014coverage}: la copertura da sola non basta.


\section{Aderenza al formato e troncamento}
% Quanto i modelli rispettano i vincoli del prompt, e cosa succede quando non lo
% fanno: la violazione del vincolo sul numero di test porta all'esaurimento del
% budget di token e quindi a file troncati, che risultano non eseguibili.
% Notare che un file troncato puo' superare il controllo sintattico se il taglio
% cade a fine funzione: sembra sano ma e' incompleto.


\section{Esempio end-to-end}
% Una funzione del dataset, il prompt costruito, il test generato dai tre
% modelli e le metriche corrispondenti, per mostrare concretamente cosa
% significa ciascuna misura.


\section{Confronto con i risultati pubblicati}
% \cite{huang2025ult} riporta, come media su 12 modelli di dimensioni maggiori e
% in gran parte specializzati sul codice, Pass@1 del 41,32% e copertura di riga
% del 45,10%. Discutere la coerenza dei risultati ottenuti e le differenze di
% protocollo che rendono i numeri non direttamente confrontabili.
